\documentclass[10pt]{article}

\usepackage[utf8]{inputenc}

\usepackage[T1]{fontenc}

\usepackage{amsmath}

\usepackage{amsfonts}

\usepackage{amssymb}

\usepackage[version=4]{mhchem}

\usepackage{stmaryrd}

\usepackage{mathrsfs}



\begin{document}

морфизма по модулю нуль (эндоморфизма $\bmod 0)$; его можно превратить в Э., выбросив из пространства подходящее множество нулевой меры. Примером Э. вероятностного пространства служит сдвиг влево в пространстве реализаций стационарной в узком смысле случайной последовательности $X_{n}, n \geqslant 0$.



Б. М. Гуревич.



OHEPГETИЧECKAЯ ЗОНА - см. Блоха закон дисперсии. ЭНЕРТЕТИЧЕСКАЯ ЩЕЛЬ - см. БКШІ-Боголюбова уравнение.



OHEPТЕТИЧЕСКОЕ УСЛОВМЕ - см. Сингулярности пространства-времени.



OHEPГИИ ИМПуЛЬСА TEHЗОР - см. Нетер теорема.



ЭНЕРГИИ ИНТЕГРАЛ - см. Динамическая система.



ЭНЕРГИИ ТУРБУЛЕНТНОСТИ УРАВНЕНИЕ - эВОЛЮционное уравнение для средней плотности кинетической энергии турбулентного течения $E=(1 / 2) \rho\left\langle u_{\alpha}^{\prime} u_{\alpha}^{\prime}\right\rangle$ (здесь $\rho-$ плотность, $u_{i}^{\prime}$ - флуктуации скорости). Для несжимаемой жидкости Э. т. у. имеет вид:



\[

\begin{aligned}

& \frac{\partial E}{\partial t}+\frac{\partial}{\partial x_{\alpha}}\left(E\left\langle u_{\alpha}\right\rangle+(1 / 2) \rho\left\langle u_{\beta}^{\prime} u_{\beta}^{\prime} u_{\alpha}^{\prime}\right\rangle+\left\langle p^{\prime} u_{\alpha}^{\prime}\right\rangle-\left\langle u_{\beta}^{\prime} \sigma_{\alpha \beta}^{\prime}\right\rangle\right)= \\

&=\rho\left\langle u_{\alpha}^{\prime} X_{\alpha}^{\prime}\right\rangle-\rho\langle\varepsilon\rangle-\rho\left\langle u_{\alpha}^{\prime} u_{\beta}^{\prime}\right\rangle \frac{\partial\left\langle u_{\alpha}\right\rangle}{\partial X_{\beta}},

\end{aligned}

\]



где $p$ - давление, $\sigma_{i j}-$ тензор вязких напряжений, $\rho X_{i}-$ плотность внешних сил, $\rho\langle\varepsilon\rangle-$ средняя удельная вязкая диссипация энергии турбулентности.



Э. т. у. в дополнение к уравнениям Рейнольдса впервые предложено А. Н. Колмогоровым [1].



Лum.: [1] Кол м о го р о в А. Н., «Известия АН СССР. Сер. физическая», 1942, т. 6, № 1-2, с. 56-8 (Избранные труды. Математика и механика, М., 1985, с. 294-96); [2] М о н и н А. С., Я гл о м А. М., Статистическая гидромеханика. Механика турбулентности, ч. 1, М., 1965.



В. П. Красицкий.



SHEPГИЯ - общая количественная мера движения и взаимодействия всех видов материи. Э. не возникает из ничего и не исчезает, она может только переходить из одной формы в другую.



В соответствии с различными формами движения материи рассматривают разные формы Э.: механическую, внутреннюю, электромагнитную, химическую, ядерную и др. Это деление до известной степени условно. Так, химич. Э. складывается из кинетич. Э. движения электронов и электрич. Э. взаимодействия электронов друг с другом и с атомными ядрами. Внутренняя Э. равна сумме кинетич. Э. хаотич. движения молекул относительно центра масс тел и потенциальная Э. взаимодействия молекул друг с другом. Э. системы однозначно зависит от параметров, характеризующих состояние системы. В случае непрерывной среды или поля вводятся понятия плотности эн е рг и и, то есть Э. в единице объема, и плотности потока эне ргии, равной произведению плотности Э. на скорость ее перемещения.



Теория относительности показала, что Э. тела \& неразрывно связана с его массой $m$ соотношением



\[

\mathscr{E}=m c^{2}

\]



Любое тело обладает $Э$.; если масса покоящегося тела $m_{0}$, то его Э. покоя



\[

\mathscr{E}_{0}=m_{0} c^{2}

\]



эта Э. может переходить в другие виды Э. при превращениях частиц (распадах, ядерных реакциях и т. п.).



Согласно классич. физике, Э. любой системы меняется непрерывно и может принимать любые значения. Квантовая теория утверждает, что Э. микрочастиц, движение к-рых происходит в огранич. объеме пространства (напр., электронов в атоме), принимает дискретный ряд значений. Так, атомы испускают электромагнитную Э. в виде дискретных порций - световых квантов, или фотонов.



678 ЭHEPTИЯ Лum.: [1] М а й е р Ю. Р., Закон сохранения и превращения энергии. Четыре исследования. 1841-1851, М.-Л., 1933; [2] Ге льмгольц Г., О сохранении силы, пер. с нем., 2 изд., М.- Л., 1934; [3] Пл а н к М., Принцип сохранения энергии, пер. с нем., М.-Л., 1938; [4] Л а уэ М., История физики, пер. с нем., М., 1956; [5] В и гн е р Е., Этюды о симметрии, пер. с англ., М., 1971. Г. Я. Мякишев. ЭНЕРГИЯ покоя - см. Релятивистская механика.



ЭНЕРГИЯ с и с те м Ы - см. Статистический ансамбль. ЭНЕРТОДОМИНАНТНОСТИ УСЛОВИЕ - см. Горизонтов топология.



ЭНСТРОФия - величина $\Omega$, равная половине среднего квадрата вихря скорости $\boldsymbol{v}$ :



\[

\Omega=\frac{1}{2}|\operatorname{rot} v|^{2} \text {. }

\]



Э. наряду с энергией является адиабатич. интегральным инвариантом в двумерной турбулентности, благодаря чему создается каскад Э. в малые масштабы по спектральному «закону минус трех», а энергии - в большие масштабы по спектральному «закону пяти третей» (см. [1]). В трехмерной турбулентности при сохраняемости энергии Э. возрастает за счет эффекта растяжения вихревых нитей. При этом энергия передается только в малые масштабы; в инерционном интервале масштабов выполняется закон Колмогорова Обухова пяти третей (см. [2]).



См. также Геофизическая турбулентность



Лum.: [1] Kraichnan R. H., «Phys. Fluids», 1967, v. 10, № 7, p. 1417-23; [2] М он и н А. С., Я глом А. М., Статистическая гидромеханика. Механика турбулентности, ч. 2, М., 1967.



А. П. Мирабель.



ЭНТАЛЬПИЯ - термодинамический потенциал, однозначная функция состояния термодинамической системы при условии, что в качестве независимых термодинамических параметров выбраны энтропия $S$ и давление $p$. Э. связана с внутренней энергией $U$ соотношением



\[

H=U+p V

\]



где $V$ - объем системы. При изобарных процессах изменение Э. равно поглощенному системой количеству теплоты $(\delta Q)_{p}:$



\[

(d H)_{p}=(\delta Q)_{p}=C_{p} d T,

\]



поэтому Э. называют те пло в ой функцие й, или те п л ос о де р жа н и е м. Э. минимальна в состоянии термодинамич. равновесия (при постоянных $p$ и $S$ ). При адиабатно-изобарных процессах в сложной системе убыль Э. равна работе немеханич. сил.



А. В. Солдатов.



ЭНТPOПИИ ЗАКОН ВОЗРАСТАНИЯ - фундаментальная гипотеза неравновесной термодинамики: в ходе спонтанной эволюции изолированной системы энтропия может лишь возрастать с течением времени в результате необратимых процессов, идущих в системе. Возрастание прекращается только тогда, когда система приходит в равновесное состояние; при этом энтропия достигает максимума. По существу, Э. з. в. есть следствие второго начала термодинамики (см. Термодинамики законы). Э. з. в. может быть доказан для нек-рых кинетич. уравнений (уравнения Больцмана и Ландау). Часто Э. з. в. называют $H$-т е о ре м о й Б о л ь цм а Ha.



Лum.: [1] Б а л е с к у Р., Равновесная и неравновесная статистическая механика, пер. с англ., т. 2, М., 1978; [2] 3 уб а ре в Д. Н., Неравновесная статистическая термодинамика, М., 1971.



ЭНТРОПИИ ПРОИЗВОДСТВО - см. Энтропия.



Д. П. Санкович. ЭНТРОПИЙНАЯ ВОЛНА - см. Магнитогидродинамические волны.



ЭНТРОПИЙНОЕ СООТНОШЕНИЕ НЕОПРЕДЕЛЕННОСТЕЙ - см. Неопределенностей соотношение.



כHTPOПИя - теоретико-информационная мера степени неопределенности случайной величины. Если $\xi \in(\Omega, \sqrt{,}, \mathrm{P})-$ дискретная случайная величина, принимающая значения $x_{1}, \ldots, x_{n}$ с распределением вероятностей $p_{1}, \ldots, p_{n}$ (где $\left.p_{i}=\mathrm{P}\left\{\xi=x_{i}\right\}\right)$, то энт ропия определяется выражением



\[

H=H\left(p_{1}, \ldots, p_{n}\right)=\sum_{i=1}^{n} p_{i} \log \left(1 / p_{i}\right)

\]



(считается, что $0 \cdot \log 0=0$ ).



М. А. Сташенко





\end{document}